\section{Proposed Deep Reinforcement Learning Architecture}
\label{sec:proposed_architecture}

To address the Resource Block Allocation (RBA) problem in Open RAN environments, we model the system as a Markov Decision Process (MDP) and propose a Deep Reinforcement Learning (DRL) agent equipped with a Graph Neural Network (GNN) backbone. The core motivation for employing GNNs is their inherent permutation equivariance and ability to process inputs of varying dimensions, enabling \textit{zero-shot generalization} to network topologies with different user densities, a critical requirement for scalable O-RAN intelligent controllers.

\subsection{State Representation as a Communication Graph}
At each Transmission Time Interval (TTI) $t$, the network state is abstracted as a directed graph $\mathcal{G}(t) = (\mathcal{V}, \mathcal{E}(t))$, where the set of vertices $\mathcal{V}$ corresponds to the $V$ active User Equipments (UEs) in the cell.

For each node $i \in \mathcal{V}$, we define a feature vector $\mathbf{x}_i(t) \in \mathbb{R}^{F}$, encapsulating local state information:
\begin{equation}
    \mathbf{x}_i(t) = \left[ \text{CQI}_i(t), Q_i(t), \lambda_i \right]^T
\end{equation}
where $\text{CQI}_i(t)$ represents the Channel Quality Indicator (acting as a proxy for direct channel gain), $Q_i(t)$ is the current buffer queue size in bits, and $\lambda_i$ denotes the target Quality of Service (QoS) requirement (e.g., CBR load profile in bytes). These vectors form the node feature matrix $\mathbf{X}(t) \in \mathbb{R}^{V \times 3}$.

The edges $\mathcal{E}(t)$ capture the spatial interference relationships between UEs. The adjacency matrix $\mathbf{A}(t) \in \mathbb{R}^{V \times V}$ contains continuous edge attributes $e_{i,j}(t)$, representing the inter-UE interference gain (CSI) from UE $i$ to UE $j$.

\subsection{Graph Attention Encoder}
To extract meaningful representations from $\mathcal{G}(t)$, we employ a Graph Attention Network (GAT) encoder. Unlike standard Convolutional Neural Networks (CNNs) or Multilayer Perceptrons (MLPs) that require fixed-size inputs, the GAT operates through a message-passing paradigm, making it structurally invariant to the number of nodes $V$.

The encoder consists of $L=2$ stacked GAT convolutional layers. For a target node $i$, the attention coefficients $\alpha_{i,j}$ with respect to its neighbor $j$ are computed by incorporating the edge attributes (interference gain), effectively allowing the network to prioritize neighbors that cause more severe co-channel interference.
The hidden representation $\mathbf{h}_i^{(l)}$ at layer $l$ is updated via:
\begin{equation}
    \mathbf{h}_i^{(l)} = \text{LayerNorm} \left( \mathbf{h}_i^{(l-1)} + \text{ELU} \left( \sum_{j \in \mathcal{N}(i) \cup \{i\}} \alpha_{i,j}^{(l)} \mathbf{W}^{(l)} \mathbf{h}_j^{(l-1)} \right) \right)
\end{equation}
where $\mathcal{N}(i)$ is the neighborhood of $i$, $\mathbf{W}^{(l)}$ is a learnable weight matrix, and residual connections are employed to mitigate over-smoothing in dense interference graphs. Multi-head attention ($H=4$ heads) is utilized to stabilize the learning process, with the outputs of the heads averaged to maintain a consistent dimensionality. The final layer outputs the node embeddings $\mathbf{H}^{(L)} \in \mathbb{R}^{V \times d_{hidden}}$ (where $d_{hidden}=64$).

\subsection{GNN-based Actor-Critic Architecture}
The extracted graph embeddings $\mathbf{H}^{(L)}$ are fed into an Actor-Critic architecture to optimize the RBA policy.

\subsubsection{Actor Network (Policy)}
The Actor is responsible for generating a probability distribution over the available UEs for each Resource Block (RB). To maintain invariance to $V$, a shared multi-layer perceptron (MLP) head processes the embedding of each node individually to produce a scalar logit:
\begin{equation}
    z_i = \text{MLP}_{\theta_{actor}}(\mathbf{h}_i^{(L)}) \quad \forall i \in \mathcal{V}
\end{equation}
A Softmax activation is then applied across the $V$ logits to form a categorical distribution $\pi_\theta(a|s)$. To allocate the $K=50$ available RBs, we independently sample from this shared distribution $K$ times, producing the final action vector $\mathbf{a}(t) \in \{0, 1, \dots, V-1\}^K$. This design ensures that the policy can seamlessly scale to any number of users.

\subsubsection{Critic Network (Value Function)}
The Critic estimates the state-value function $V(s)$. To summarize the global state of the cell regardless of the number of active UEs, a global mean pooling operation aggregates the node embeddings into a single graph-level representation, which is subsequently processed by an MLP to output a scalar value:
\begin{equation}
    V_\phi(s) = \text{MLP}_{\phi_{critic}} \left( \frac{1}{V} \sum_{i \in \mathcal{V}} \mathbf{h}_i^{(L)} \right)
\end{equation}

\subsection{Action Space and Reward Formulation}
To fully define the underlying Markov Decision Process (MDP), we specify the action space and the reward signal that guides the agent's learning process.

\subsubsection{Multi-Discrete Action Space}
The Open RAN Medium Access Control (MAC) scheduler must assign $K$ available Resource Blocks (RBs) to the $V$ active UEs at each TTI. Instead of a highly combinatorial action space, which scales poorly, we formulate a multi-discrete action space $\mathcal{A}$. An action $\mathbf{a}(t) = [a_1, a_2, \dots, a_K]$ is defined such that each element $a_k \in \{0, 1, \dots, V-1\}$ indicates the UE index assigned to the $k$-th RB. This allows the Actor network to sample $K$ independent categorical decisions, greatly simplifying the policy parameterization while maintaining full allocation flexibility.

\subsubsection{Reward Function for URLLC}
The objective of the DRL agent is to maximize overall system spectral efficiency while strictly minimizing the queuing delay, aligning with the stringent requirements of Ultra-Reliable Low-Latency Communication (URLLC) services. The reward function $r(t)$ at TTI $t$ is formulated as the total cell throughput penalized by the aggregate remaining queue sizes:
\begin{equation}
    r(t) = \sum_{i \in \mathcal{V}} C_i(t) - \beta \sum_{i \in \mathcal{V}} Q_i(t)
\end{equation}
where $C_i(t)$ is the realized throughput for UE $i$ (bounded by its current buffer demand and the Shannon capacity of its allocated RBs), $Q_i(t)$ is the residual queue length in bits at the end of the TTI, and $\beta$ is a tunable penalty weight (e.g., $10^{-4}$). This composite reward signal naturally guides the PPO algorithm to learn a \textit{Proportional Fair} scheduling strategy that prevents buffer overflows, effectively reducing the tail end (e.g., $95^{th}$ percentile) of the packet delay distribution.

\subsection{Baseline: MLP Actor-Critic}
For comparative analysis, we define a baseline MLP-based Actor-Critic agent. This model flattens the node features $\mathbf{X}(t)$ and the upper triangle of the adjacency matrix $\mathbf{A}(t)$ into a single fixed-dimensional vector. This flattened vector is passed through fully-connected layers to output the policy and value estimates. By design, the input dimensionality of the MLP is strictly tied to the number of UEs used during training ($V_{train}$), structurally prohibiting zero-shot inference in scenarios where $V_{test} \neq V_{train}$. This baseline serves to empirically demonstrate the necessity of the GNN architecture for topological generalization.

\subsection{Proximal Policy Optimization (PPO)}
Both agents are trained using Proximal Policy Optimization (PPO), a state-of-the-art policy gradient method. PPO optimizes a surrogate objective function that prevents destructively large policy updates by clipping the probability ratio $r_t(\theta) = \frac{\pi_\theta(a_t|s_t)}{\pi_{\theta_{old}}(a_t|s_t)}$. 
The overall loss function minimized during training over iterations with mini-batches is:
\begin{equation}
    \mathcal{L}^{CLIP+VF+S}(\theta) = \mathbb{E}_t \left[ \mathcal{L}^{CLIP}_t(\theta) - c_1 \mathcal{L}^{VF}_t(\theta) + c_2 S[\pi_\theta](s_t) \right]
\end{equation}
where $\mathcal{L}^{CLIP}_t(\theta) = \min(r_t(\theta)\hat{A}_t, \text{clip}(r_t(\theta), 1-\epsilon, 1+\epsilon)\hat{A}_t)$, $\hat{A}_t$ is the Generalized Advantage Estimation (GAE- $\lambda$) computed using the Critic's evaluations from the rollout buffer, $\mathcal{L}^{VF}_t$ is the value function squared-error loss, and $S[\pi_\theta]$ is an entropy bonus to encourage exploration and prevent premature convergence to suboptimal deterministic policies.
